\section{What the Composer Checks, and What It Does Not}
\label{sec:ch15-what-the-composer-checks}
\index{composition!authorization directives|(}%

Everything so far has been one scope on one field, working. This section is the
other half of the job, which is finding out what the composer will let past
when the rule is wrong. Ten cases, in \code{scripts/auth-cases.mjs}: four read
the committed schemas as they stand, and six edit them the way
chapter~\ref{ch:composition} and chapter~\ref{ch:hard-modeling-problems} did
before this. Three of the six compose cleanly and should not.

\subsection*{The import list is decoration}

\index{link@\code{@link}!import list ignored for authorization directives}%
Mosaic's schemas each carry an \code{@link} naming the federation version and
importing the directives they use. Take \code{"@requiresScopes"} out of
Accounts' import list, leave the directive on the field, and the graph composes
with the same authorization configuration as before.

That is not the composer being relaxed about unknown directives. The control
case is one line away:

\begin{minted}[fontsize=\footnotesize,breaklines]{text}
The subgraph "accounts" could not be federated for the following reason:
The directive "@totallyMadeUpDirective" declared on coordinates
"Customer.displayName" is not defined in the schema.
\end{minted}

The two authorization directives are registered unconditionally. The composer
at 0.63.2 has a set with exactly two members in it, and the import list is
never consulted about either. That is harmless in itself, and worth knowing
because it means a subgraph's \code{@link} is not a reliable inventory of what
that subgraph is doing.

\subsection*{An empty list is not an empty rule}

\index{requiresScopes@\code{@requiresScopes}!empty scopes composes silently}%
Change the argument to \code{scopes: []}. It composes with no error and no
warning, and the directive is gone: not weakened, not defaulted, absent from
the composed client schema and absent from \code{fieldConfigurations}. The
field is public again and the graph is happy.

This is the one that would actually happen. Somebody writes a list, a
refactoring empties it, a template renders nothing into it, and the protection
evaporates with every check in the toolchain green. There is no diff in the
router's behaviour to notice, because the router was never told.

\subsection*{One nesting level short crashes it}

\index{composition!TypeError on a flat scopes list}%
The nesting is the part everybody gets wrong, so it is worth knowing what
getting it wrong looks like. \code{scopes: ["read:pii"]}, one level short of
the \code{[[Scope!]!]!} the directive declares, does not produce a validation
message. It throws:

\begin{minted}[fontsize=\footnotesize,breaklines]{text}
TypeError: scopes.values is not iterable
    at NormalizationFactory.extractRequiredScopes
\end{minted}

wrapped in the composer's generic \enquote{upgrade to the latest version, and
if it persists please open an issue} box. Chapter~\ref{ch:hard-modeling-problems}
met the same class of failure with an entity interface missing a key, and the
same thing is true here: a stack trace is not a diagnosis, and the box asks you
to report a bug in your own schema.

Omit the argument entirely and the composer behaves properly, which shows the
validation exists and only that one shape misses it:

\begin{minted}[fontsize=\scriptsize,breaklines]{text}
The definition for "@requiresScopes" defines the following 1 required argument:
"scopes". However, no arguments are defined on this instance.
\end{minted}

\subsection*{\code{@policy} is not available here}

\index{policy@\code{@policy}!not implemented by this composer}%
HotChocolate ships \code{[Policy]} and emits \code{@policy}, against the
federation version its own source marks the directive as belonging to. This
stack cannot compose it:

\begin{minted}[fontsize=\footnotesize,breaklines]{text}
The subgraph "accounts" could not be federated for the following reason:
The directive "@policy" declared on coordinates "Customer.displayName"
is not defined in the schema.
\end{minted}

with or without an \code{@link} import naming it, because the composer's
registry of authorization directives has two entries and consults no import
list to learn a third. The router agrees: there is no policy evaluation
anywhere in its authorizer.

So a HotChocolate subgraph author reading their own attribute list will find
three federation authorization attributes, of which two work here and the third
stops the graph from composing at all. That is worth stating plainly because
the failure names the directive and not the reason, and a reasonable person
would conclude they had spelled something wrong.

\subsection*{Two subgraphs, two scopes, and a field that reads the other way}

\index{requiresScopes@\code{@requiresScopes}!merging across subgraphs}%
\code{Money.amount} is \code{@shareable} and declared in both Pricing and
Ordering. Give it one scope in the first and a different one in the second,
and the graph composes with no error and no warning. What comes out is this:

\begin{minted}[fontsize=\footnotesize,breaklines]{text}
"requiredOrScopes":     [ { "requiredAndScopes": ["read:orders", "read:price"] } ]
"requiredOrScopesByOr": [ { "requiredAndScopes": ["read:price"] },
                          { "requiredAndScopes": ["read:orders"] } ]
\end{minted}

Read those carefully, because they disagree. The first says a caller needs both
scopes. The second says either one will do. They sit next to each other in one
JSON object, written by one composer, from one pair of schemas, and nothing
anywhere flags that the graph has two answers.

Which one is real? The router reads the first of the two and nothing else. The
second appears in the router's Go source only as a generated protobuf
accessor, and is never called. So the enforced rule is the strict one, and the
field whose name reads like the obvious home for an \enquote{OR of ANDs}
scheme is dead weight in the file.

Be careful with what that means, because \enquote{it merges to AND} sounds
like a design choice and I am not sure it is one. What I measured is that the two subgraphs' requirements combine into a
single AND-group, and that a caller therefore needs a scope neither service
asked for on its own. If your mental model was that each service guards its own
contribution, that model is wrong here, and the composer will not tell you.

\subsection*{And the field this chapter did not fix}

\index{node@\code{Query.node}!authorising it}%
Chapter~\ref{ch:hard-modeling-problems} handed this chapter \code{Query.node}
as a named risk and I am handing it on, so it is worth saying exactly why
rather than leaving it in a list.

\code{node} takes one identifier and returns anything in the graph. Its type is
known only after the identifier has been decoded, which happens inside the
service that owns the type, which is after any rule the router could apply.
So the two available rules are both wrong. A scope on \code{node} itself
governs every type at once, which either blocks identifiers a caller is
entitled to or admits ones they are not. A scope on each concrete type governs
the type wherever it appears, including on the paths where it is reached
normally and was fine.

The shape of an answer is probably that the node service refuses types rather
than the router refusing the field, since it is the only process that knows
which type an identifier named, and it owns nothing else. That is a resolver
making an authorisation decision with a claim it was handed by the header rule
of section~\ref{sec:ch15-the-header-nobody-forwards}, and it is invisible to
the composer, which is exactly the arrangement this chapter has spent six
sections warning about. I do not have a version of it I would defend, so
\code{node} is unprotected at tag \code{ch15} and
chapter~\ref{ch:inside-the-router-and-the-composer} inherits the question with
the composer's own rules in front of it.

\medskip

What this chapter did close is the gap it opened with: Mosaic has a rule the
router enforces, a rule a service enforces, and a written argument for which
kind of rule belongs where. The lab is where you take each of them apart.

\index{composition!authorization directives|)}%
